% New citation records introduced in this section are supplied in
% Chapter2/chapter2_literature_review_additions.bib. The thesis-level
% bibliography, which is not present in this repository snapshot, must include
% that database when the master document is restored.

\section{Literature Review}\label{ch2:sec:lr}

This chapter lies at the intersection of several literatures that study different objects and support different empirical claims. 

Research on central-bank communication asks what information is contained in observed policy documents and how markets respond when those documents are released. Research on monetary-policy forecasting asks whether an action can be predicted from information available before a meeting. 

Work on data-to-text generation instead asks whether a model can transform a supplied record into a fluent and faithful text, while the post-training literature studies how supervised or reward-based objectives alter model behaviour. 

These questions are related, but their estimands are not interchangeable. In particular, a synthetic paragraph that resembles an FOMC Minute is neither an observed communication shock nor evidence that the paragraph affected historical asset prices. The review therefore separates the relevant strands before locating the two parallel tasks studied in this chapter.


\subsection{Synthetic Economic Data and Controlled Data-to-Text Generation}
\label{ch2:sec:lr:synthetic}

Synthetic data are generated observations intended to reproduce selected properties of an empirical data-generating process without being literal copies of the original observations. In economics and finance, their uses include simulation, robustness analysis, privacy protection, rare-event augmentation, and stress testing. These uses impose different validity requirements. A sample that is useful for testing a forecasting pipeline need not reproduce every feature of the source distribution, while data intended for substantive counterfactual inference require a much stronger structural argument. Financial applications illustrate the distinction: synthetic limit-order books and price paths can be useful for controlled robustness exercises even when they are not treated as newly observed market histories (\cite{cao2022synthetic}; \cite{carvajal2022synthetic}). Likewise, pretrained language models can generate additional labelled examples or task demonstrations (\cite{schick2021generating}), but the utility of those examples depends on the fidelity of the generator, the selection rule, and the downstream task.

One particularly relevant methodological tradition for the Minutes branch is
consequently
\emph{controlled data-to-text generation}, rather than unconstrained document
simulation. Natural-language generation research distinguishes content
selection, document planning, surface realization, and evaluation
(\cite{gatt2018survey}). Neural data-to-document systems can produce fluent
prose while omitting source records, introducing unsupported relations, or
misstating values; \citet{wiseman2017challenges} therefore evaluate both
linguistic quality and record-level content fidelity. This distinction is
especially important for central-bank prose, where a misplaced sign, date, or
rate can reverse the economic interpretation even if the paragraph remains
stylistically convincing. Accuracy protocols for data-to-text systems
accordingly treat factual error annotation and metric validation as separate
from fluency assessment (\cite{thomson2020gold}).

This literature narrows the claim made in the present chapter. Model chk-2
does not generate a complete institutional record from an unconstrained
counterfactual economy. It rewrites a supplied macro-financial analysis into
one Minutes-style paragraph under a fixed output contract. The relevant
question is thus whether task-specific training improves preservation of the
supplied content and alignment to a fixed source-grounded target. Style
similarity is useful only after delivery and source-preservation requirements
are checked. Moreover, because the evaluation targets are synthetic rather
than official FOMC Minutes, agreement with them establishes target alignment,
not equivalence to the Committee's historical language or deliberative
process.

\subsection{Central-Bank Communication, Information, and Asset Prices}
\label{ch2:sec:lr:communication}

Central-bank communication is part of the monetary-policy transmission
mechanism because it can shape expectations about both the current instrument
and its future path. The survey by \citet{blinder2008central} concludes that
communication can move financial markets and improve the predictability of
policy, while also emphasizing that communication strategies and their
effects vary across institutions and instruments. High-frequency evidence
reinforces the need to separate an action from its accompanying message.
\citet{gurkaynak2005actions} show that U.S. announcement responses cannot be
summarized by a single target-rate surprise: a distinct path factor, closely
associated with FOMC statements, helps explain longer-maturity yields. At the
same time, the central bank may reveal information about economic fundamentals
that private agents do not possess (\cite{romer2000federal}). The
\emph{information effect} identified by \citet{nakamura2018information}
therefore warns against reading every announcement-window price movement as a
pure policy shock.

Communication content is itself multidimensional. \citet{hansen2016shocking}
separate language about current economic conditions from forward guidance and
show that the dimensions have different yield effects. A scalar sentiment
index can therefore be a useful summary without being a sufficient statistic
for the economic content of a policy document.

Minutes are also an institutional record of deliberation rather than simply a
longer policy statement. Using the 1993 transparency change and computational
linguistics, \citet{hansen2018transparency} show that publication rules can
alter deliberative communication itself. This result matters for synthetic
generation: matching the surface form of published Minutes does not recreate
the incentives, information, or bargaining process that produced the official
document. It also cautions against treating a generated paragraph as a
measurement-free proxy for an unobserved institutional state.

Nor are Minutes verbatim transcripts. \citet{apel2022disclosure} compare FOMC
Minutes with the underlying transcripts and show that the shortened, edited
Minutes can omit committee, chair, and disagreement signals that retain
predictive content for future policy. Institutional style is therefore not a
lossless encoding of the meeting. A model trained to reproduce the published
genre can at most learn the disclosure convention represented in its targets;
it cannot be assumed to recover omitted deliberative information.

Within the empirical text-as-data literature, the unit of analysis and the
release clock determine the interpretation. Event studies examine asset-price
movements when statements or Minutes become public. Existing work reports
effects on exchange rates, yields, returns, or volatility around Federal
Reserve communications (\cite{rosa2011high}; \cite{rosa2013financial};
\cite{jubinski2013fomc}). Other studies convert observed Minutes into
sentiment or topic measures and assess their predictive content.
\citet{huang2021economic}, for example, construct mandate-specific indicators
and study forecasts of macroeconomic variables, whereas
\citet{tadle2022fomc} relates policy-document sentiment to fed funds futures
and foreign-exchange valuations. These designs study realized official text.
They do not validate a generated document merely because a classifier assigns
it a similar tone.

Recent work brings LLM measurement directly into this literature.
\citet{chen2026decoding} use an LLM-based framework to measure incremental
linguistic information in actual Minutes relative to same-meeting statements
and combine that measure with a one-minute Treasury-futures event study. This
is a particularly close observed-document application among the studies
reviewed here, but its identification comes from the staged release of two
official texts. The
present Minutes branch instead evaluates synthetic rewriting of a supplied
analysis and has no historical release event of its own.

Text scores are also generated covariates rather than error-free observables.
\citet{christensen2026valid} show that ignoring error in AI- or
machine-learning-generated regressors can bias estimation and invalidate
inference, including in an FOMC-sentiment application. Accordingly, agreement
between coefficient signs across generated-text artifacts is weaker evidence
than a validated, error-corrected estimate based on observed communication.

The distinction is decisive for the historical sentiment--Treasury exercise
in this chapter. The synthetic paragraphs were not released to investors at
the historical event time. A regression using their sentiment scores can
compare reduced-form coefficient patterns across text artifacts, but it cannot
identify a market response caused by the synthetic text. It also cannot, by
itself, establish factual grounding: two documents can induce similar
sentiment scores while differing in economically important claims. The
exercise is therefore an association diagnostic motivated by the observed-text
literature, not a synthetic event study.

\subsection{Real-Time Information Sets and Monetary-Policy Prediction}
\label{ch2:sec:lr:realtime}

Monetary-policy prediction requires a stricter timing convention than
retrospective text classification. Revised macroeconomic series can contain
information that policymakers and market participants did not possess when a
decision was made. \citet{croushore2001realtime} formalize the use of data
vintages, and \citet{orphanides2001monetary} demonstrates that policy-rule
recommendations and estimated historical reaction functions can change
materially when real-time rather than revised data are used. A historically
accurate model can therefore be information-infeasible if its inputs include
later releases, revisions, or the target meeting's outcome.

The action being predicted must also be defined. Econometric work often
separates whether the target rate changes from the conditional size of the
change (\cite{hamiltonjorda2002target}). Research using Minutes to forecast
subsequent federal funds rate changes, such as \citet{jung2016minutes}, asks
whether released official documents revise expectations about a later
meeting. That problem differs from the decision branch in this chapter, which
maps a target-neutral pre-meeting analysis directly to the action at the same
meeting. Here, \emph{target-neutral} means that explicit target actions and
post-meeting observations are excluded; it does not prove that meeting
identity is impossible to infer from the remaining historical context.

Established forecasting designs make these timing requirements operational.
\citet{hayo2010federal} augment a real-time, forward-looking Taylor-rule
ordered probit with Federal Reserve communications, while
\citet{vandendhauwe2013bayesian} combine real-time macro-financial predictors,
dynamic ordered probit models, and Bayesian variable selection in an
out-of-sample decision exercise. A particularly close text-informed action
comparator is \citet{aruoba2024identifying}, who use NLP on staff documents prepared
before policy decisions to approximate the Federal Reserve's information set
and predict target-rate changes. Their prediction is an intermediate step in
monetary-shock identification, and the internal staff documents were not
available to external forecasters at the decision time. Market-based studies
provide another benchmark: fed funds futures isolate anticipated and surprise
components (\cite{kuttner2001monetary}), while the yield curve and meeting
calendar also contain information about policy decisions
(\cite{piazzesi2005bond}).

These findings motivate three design requirements. First, inputs must be bound
to a point-in-time cutoff rather than reconstructed from final-vintage data.
Second, a prospective forecasting study should separate train, validation,
and test meetings chronologically so that later meetings do not teach the
model how earlier information maps to a future policy regime. Third,
evaluation must reflect the ordered but imbalanced cut--hold--hike action
space. Raw accuracy can be dominated by the majority hold class; balanced
accuracy, macro F1, class-wise recall, exact action, and output delivery
consequently answer different questions. The current release is date ordered,
but its 13-meeting test has already been opened and is class sparse; it is
therefore a training-held-out historical check rather than a new prospective
validation. In addition, the pre-2009 supplement enforces its cutoff relative
to the workbook's stored date, which is the meeting end date for 25 of 109
meetings, not uniformly the official start date. These qualifications prevent
the exercise from being interpreted as a structural policy rule or a causal
estimate of how the FOMC would respond in an unseen counterfactual economy.

\subsection{Domain Adaptation, Synthetic Supervision, and Model Lineage}
\label{ch2:sec:lr:adaptation}

Broadly pretrained language models need not encode the terminology, discourse
structure, or evidentiary discipline of a specialized task. Domain- and
task-adaptive training can improve downstream performance even after general
pretraining (\cite{gururangan2020dont}). Parameter-efficient methods such as
LoRA and QLoRA make such adaptation feasible by updating low-rank components
rather than every base-model parameter (\cite{hu2021lora};
\cite{dettmers2024qlora}). These methods explain how a comparatively small
institutional corpus can alter output behaviour, but they do not guarantee
that the resulting model is faithful or that an apparent improvement is due
to the intended capability.

The provenance of the supervision therefore determines what a checkpoint
contrast can identify. Knowledge
distillation transfers behaviour from a teacher to a smaller or otherwise
different student (\cite{hinton2015distilling}); language-model-generated
training pairs extend this idea to synthetic supervision. Self-generated
instructions and rationales can provide useful task supervision
(\cite{wang2023selfinstruct}; \cite{hsieh2023distilling}), but imitation may
transfer a teacher's response conventions more readily than its underlying
factual capabilities (\cite{gudibande2024false}). Such targets can scale
annotation and impose a consistent response contract, but they can also
transmit teacher omissions, factual errors, stylistic artefacts, and prompt
dependencies. A student--teacher similarity score is consequently not an
independent truth criterion. The teacher prompt, generation settings, frozen
outputs, split discipline, and post-generation audits must be treated as part
of the data provenance.

This chapter's shared-parent design is motivated by that concern. Model chk-1
provides the same analytical parent for two independent branches, so the main
Minutes contrast isolates the incremental association of Minutes-rewriting SFT
with output metrics relative to Model chk-1. The decision branch likewise
compares an SFT state with its own post-GRPO descendant. These are
checkpoint-specific contrasts, not generic estimates of \emph{fine-tuning} or
\emph{reinforcement learning}. They also do not establish a serial
indicators-to-Minutes-to-decision pipeline, because Model chk-3 never consumes
Model chk-2 outputs.

\subsection{Faithfulness, Evaluation, and Input-Sensitivity Diagnostics}
\label{ch2:sec:lr:evaluation}

The generation literature distinguishes \emph{faithfulness} to a supplied
source from \emph{factuality} with respect to the external world. A statement
may be entailed by an erroneous source, or may be externally true but
unsupported by the provided input. Human evaluations of abstractive systems
show that fluent outputs can still contain substantial unsupported content
(\cite{maynez2020faithfulness}), and broader reviews document the persistence
of hallucination under modern language modelling (\cite{huang2025survey}).
For this chapter, at least four constructs must therefore remain separate:
source preservation, external factuality, target-text similarity, and
institutional style or delivery. None is an automatic substitute for the
others.

Automatic metrics operationalize only parts of these constructs. N-gram
scores measure surface overlap, while embedding-based measures such as
BERTScore reward contextual similarity (\cite{zhang2019bertscore}). They can
recognize valid paraphrases, but a high semantic score does not ensure that
every number, date, direction, or attribution is correct. Conversely, one
lexically different but faithful realization may score below a stylistically
similar text containing a critical numerical error. The data-to-text
evaluation literature therefore recommends validating automatic metrics
against explicit accuracy judgements (\cite{thomson2020gold}). The present
chapter implements this construct separation by reporting hard structure,
numeric, date, and degeneration gates alongside reliability-adjusted semantic
scores. Zeroing semantic scores after a hard-gate failure encodes a declared
reliability rule; it does not transform the semantic metric itself into a
factuality test.

Reference provenance further limits interpretation. The external
Minutes-rewriting targets used here are deterministic, synthetic, and
source-grounded, not official Minutes. They permit a controlled paired
comparison among checkpoints, but agreement may reflect shared target design
choices. The evaluation release is independent of chapter-specific
fine-tuning and checkpoint selection, yet this chapter cannot establish that
its historical subject matter was absent from the pretrained base model. This
is a contamination limitation, not a reason to relabel the release as part of
chapter-specific training (\cite{magar2022contamination};
\cite{jacovi2023stop}). Similarly, the small diagnostic used to choose a
checkpoint is selection evidence rather than an unbiased population estimate,
because optimizing and evaluating on the same cases induces selection bias
(\cite{cawley2010overfitting}).

Repeated stochastic decoding is consequently used to estimate performance
over generated samples rather than one favourable completion, with common
meeting--replicate seeds preserving the paired comparison. Paired resampling
is well suited to system comparisons in NLG (\cite{deutsch2021statistical}),
and the broader machine-learning literature shows that reported performance
can vary with multiple stochastic sources (\cite{bouthillier2021accounting}).
The present hierarchy covers meeting and decoding variation; it does not cover
independent retraining variation because each stage contributes one trained
artifact. Meeting-level hierarchical resampling and family-wise multiplicity
correction reflect the nested design and multiple reported endpoints,
consistent with the requirement that NLP significance procedures match the
sampling unit and experimental comparison (\cite{dror2018hitchhiker}).

The canonical Core8 intervention would ask a different question: how target
alignment changes when one complete input block is deleted or replaced by
token-matched neutral text while a deterministic synthetic source-grounded
target remains fixed. The surviving evidence in this chapter instead
comprises two legacy deletion aggregates, one referenced to the model's
unmasked output and one to actual FOMC Minutes, whose paired seed- and
row-level provenance could not be reconstructed. Those archived targets define
different estimands from the canonical synthetic-target contract; the
neutral-text comparison and canonical paired run remain to be completed.
Shapley-value methods define feature contributions by
averaging marginal changes across coalitions (\cite{shapley1953value};
\cite{lundberg2017unified}). A single leave-one-block contrast does not perform
that averaging and can be sensitive to an out-of-distribution mask. The legacy
aggregate should therefore be interpreted only as a local prompt-sensitivity
diagnostic, not as a causal decomposition of the generated text, proof that
the model used the indicator in a human-like way, or a full Shapley
attribution. Keeping this boundary explicit connects the proposed
perturbation design to the literature without overstating what the surviving
evidence identifies.

\subsection{Reward-Based Post-Training and Structured Policy Decisions}
\label{ch2:sec:lr:decision}

Reinforcement learning can optimize sequence-level properties that are awkward
to express through token-level supervised likelihood. Classical RLHF pipelines
may learn a reward model from human rankings and optimize it with Proximal
Policy Optimization (\cite{ouyang2022training};
\cite{schulman2017proximal}). The present decision branch is different: it
uses a deterministic action-and-format reward rather than human preferences or
a learned preference model. Group Relative Policy Optimization (GRPO)
estimates relative advantages from multiple completions for the same prompt
without a separate critic (\cite{shao2024deepseekmath}). DAPO subsequently
introduced changes to clipping, sampling, loss aggregation, and handling of
long responses for large-scale reasoning-model training
(\cite{yu2025dapo}). These methods were developed primarily in settings with
verifiable mathematical or rule-based outcomes. Their use for a small,
imbalanced monetary-policy task remains an empirical transfer, not a guarantee
of improved policy reasoning.

The reward definition is central to that transfer. A deterministic reward for
JSON delivery, action direction, magnitude proximity, and exact match can make
the response contract auditable. It is nevertheless a proxy for the broader
construct of decision quality. Optimizing an imperfect proxy can increase the
proxy while reducing performance under an external criterion, a pattern
studied directly in work on reward-model overoptimization
(\cite{gao2023reward}). Format compliance and economic classification should
therefore be reported separately, component weights should not be presented as
theory-derived without evidence, and checkpoint selection should rely on
held-out outcome metrics rather than training reward alone.

Generated rationales create an additional interpretive risk. Chain-of-thought
text can be plausible while failing to reveal the causal basis of a model's
answer (\cite{turpin2023unfaithful}). In the current decision reward, a
nonempty reasoning prefix satisfies a delivery condition, but the reward does
not independently score the rationale's factual grounding, completeness, or
economic coherence. The chapter accordingly evaluates the terminal action
and native delivery contract and does not treat the emitted rationale as
evidence of interpretable or institutionally valid reasoning. This literature
also explains why a higher post-GRPO delivery rate accompanied by lower
decision metrics is a format--decision trade-off rather than an overall GRPO
improvement.

\subsection{Selected Closely Related Work and the Research Gap}
\label{ch2:sec:lr:gap}

Table~\ref{tab:ch2:closest_literature} compares selected closely related
research designs
along the dimensions that determine what each can identify. No single row is
a direct predecessor to the entire chapter. Observed-communication studies
have institutional authenticity and market timing but do not generate
controlled counterfactual text. Generic data-to-text studies provide methods
for conditional generation and fidelity assessment but do not impose a
real-time monetary-policy information set. Policy-choice models produce
actions from quantitative information but do not test institutional rewriting.
Finally, reward-based LLM training can enforce verifiable output contracts,
but an optimized proxy need not improve external policy classification.

\begin{table}[htbp]
    \centering
    \footnotesize
    \setlength{\tabcolsep}{3pt}
    \renewcommand{\arraystretch}{1.18}
    \begin{tabularx}{\textwidth}{@{}
        >{\raggedright\arraybackslash}p{0.15\textwidth}
        >{\raggedright\arraybackslash}p{0.23\textwidth}
        >{\raggedright\arraybackslash}p{0.17\textwidth}
        >{\raggedright\arraybackslash}p{0.18\textwidth}
        >{\raggedright\arraybackslash}X@{}}
        \toprule
        \textbf{Study or design} & \textbf{Task and information timing} &
        \textbf{Method} & \textbf{Evidence supported} &
        \textbf{Remaining gap relative to this chapter} \\
        \midrule
        Central-bank announcements
        (\cite{gurkaynak2005actions}; \cite{nakamura2018information})
        & Official actions and communication observed in an announcement
        window
        & High-frequency event-study and shock-identification designs
        & Asset-price responses and policy- versus information-shock
        distinctions
        & Studies realized public communication, not synthetic
        source-conditioned text. \\

        FOMC Minutes disclosure and text analysis
        (\cite{apel2022disclosure}; \cite{huang2021economic};
        \cite{tadle2022fomc})
        & Official Minutes and, for Apel et al., transcripts observed after
        the meeting
        & Dictionaries, text mining, and predictive regressions
        & Disclosure loss, tone, forecasts, and predictive content
        & Establishes that actual Minutes are edited and potentially
        incomplete; does not validate a generator. \\

        LLM analysis of observed FOMC text
        (\cite{dunn2024using}; \cite{chen2026decoding})
        & Released Statements or Minutes at or after publication
        & LLM classification or measurement combined with validation or an
        event study
        & Topic-label agreement or Treasury-futures response
        & Measures authentic documents; does not test source-conditioned
        generation or pre-meeting action prediction. \\

        Text-informed policy-action models
        (\cite{hayo2010federal}; \cite{vandendhauwe2013bayesian};
        \cite{aruoba2024identifying})
        & Real-time variables, communication, or pre-decision internal staff
        text
        & Ordered probit, Bayesian selection, NLP, and machine learning
        & Out-of-sample action prediction or shock diagnostics
        & Does not produce Minutes-style prose; the information set and
        estimand differ from this chapter. \\

        Controlled data-to-text generation
        (\cite{wiseman2017challenges}; \cite{thomson2020gold})
        & Structured records or source documents available at generation
        & Neural conditional generation with content planning or copying
        & Fluency, record coverage, and human factual-accuracy assessment
        & Establishes fidelity tools in other domains, not an FOMC-specific
        information contract. \\

        Domain adaptation and synthetic supervision
        (\cite{gururangan2020dont}; \cite{hinton2015distilling})
        & Domain corpora or teacher outputs whose timing depends on provenance
        & Continued pretraining, SFT, or distillation
        & Changes in downstream task performance
        & Gains can inherit teacher error and do not by themselves establish
        grounding. \\

        GRPO and structured-output optimization
        (\cite{shao2024deepseekmath}; \cite{yu2025dapo})
        & Prompts paired with verifiable targets or proxy rewards
        & Group-relative policy optimization
        & Reward, exact-answer, and format metrics
        & Developed in more readily verifiable domains; proxy improvement may
        diverge from policy accuracy. \\

        \textbf{This chapter}
        & Minutes branch: supplied Model chk-1 analysis on an external
        pre-2009 evaluation. Decision branch: target-neutral pre-meeting brief
        in a date-ordered release, subject to the 25/109 supplement-date caveat
        and an already-opened, class-sparse test
        & Shared analytical parent with separate rewriting-SFT and
        decision-SFT--GRPO branches
        & Hard fidelity gates, semantic alignment, bounded perturbation
        evidence, and training-held-out action metrics
        & Does not validate full-document institutional simulation, causal
        market impact, rationale faithfulness, or a serial
        Minutes-to-decision pipeline. \\
        \bottomrule
    \end{tabularx}
    \caption{Selected Closely Related Research Designs and the Remaining
    Identification Gap}
    \label{tab:ch2:closest_literature}
\end{table}

The resulting gap is narrower, and more testable, than a claim to simulate the
FOMC. Prior work largely either measures observed central-bank text or develops
general-purpose conditional generation and post-training methods. This chapter
tests two conditional mappings from supplied analyses: whether
Minutes-specific SFT improves source-preserving institutional rewriting, and
whether a separate decision-specific GRPO stage changes structured delivery
without degrading held-out policy-action quality. The first link is evaluated
with hard fidelity gates, stochastic semantic comparisons, and bounded
sensitivity diagnostics. The second is evaluated as a small historical
classification exercise in which delivery and decision quality remain
distinct. It does not treat the mappings from analysis to Minutes-style prose
and, separately, from a pre-meeting brief to a policy action as a serial causal
system.

Taken together, these literatures motivate five bounded empirical questions
rather than one general claim of institutional simulation. RQ1 separates hard
source preservation from semantic alignment in the Model chk-2 contrast. RQ2
frames Core8 deletion and neutralization as local input-sensitivity tests, not
as causal or Shapley decompositions; only the legacy deletion aggregate is
currently available. RQ3 specifies a future direct sentiment-score closeness
diagnostic for which no active result is reported. RQ4 asks what the surviving
historical sentiment--Treasury record documents, while recognizing that its
provenance is insufficient for a verified cross-artifact difference test. RQ5
evaluates delivery and action quality on the independent Model chk-3 branch
and requires the two outcomes to be reported separately. This mapping assigns
each question an evidentiary standard drawn from the relevant literature and
prevents a positive result on one construct from being used as evidence for
another.

Among the studies reviewed here, \citet{chen2026decoding} provide a
particularly close comparator on the Minutes side, but they measure
information in actual released documents rather than generate
source-conditioned prose. \citet{aruoba2024identifying} provide a particularly
close comparator on the action side, but their prediction uses internal staff
documents and serves monetary-shock construction rather than evaluation of an
externally deployable classifier. The intersection addressed
here---controlled institutional rewriting and a separately evaluated,
date-ordered action branch with the disclosed supplement-date
qualification---therefore remains comparatively underexplored.
